% !TEX root = ../main.tex
Il file \textit{stats\_utils.py} raccoglie le funzioni statistiche di base usate nel package. Le funzioni lavorano su array monodimensionali finiti; quando sono presenti i pesi, questi devono avere la stessa forma dei dati ed essere strettamente positivi.

Le funzioni disponibili sono:
\begin{itemize}
    \item \texttt{median(x)}
    \item \texttt{weighted\_mean(x, w=None)}
    \item \texttt{variance(x, w=None, ddof=0)}
    \item \texttt{covariance(x, y, w=None)}
    \item \texttt{standard\_deviation(x, w=None, ddof=0)}
\end{itemize}

\begin{quote}
\textbf{Scelta di default:} il package usa statistiche descrittive di popolazione. Per questo \texttt{ddof=0} è il default coerente in tutto \texttt{mespy}. Se serve la correzione di Bessel, passare \texttt{ddof=1} in modo esplicito.
\end{quote}

\section{Validazione comune}
\begin{apidesc}
  \item[\textbf{Contratto}]
    Tutte le funzioni statistiche convertono gli input in array \texttt{float64} monodimensionali. Gli array vuoti o con valori non finiti (\texttt{NaN}, \texttt{inf}) vengono rifiutati con \texttt{ValueError}.
  \item[\textbf{Pesi}]
    Nei casi pesati, \texttt{w} deve avere la stessa forma dei dati, contenere solo valori strettamente positivi e avere somma strettamente positiva.
  \item[\textbf{Obiettivo}]
    Evitare warning di NumPy e risultati \texttt{nan} silenziosi nell'uso didattico.
\end{apidesc}

\section{median(x)}
\label{sec:median}
\begin{apidesc}
  \item[\textbf{Scopo}]
    Calcolare la mediana di un array di dati numerici.
  \item[\textbf{Restituisce}]
    \texttt{float}: mediana di \texttt{x}.
  \item[\textbf{Eccezioni}]
    \texttt{ValueError} se \texttt{x} è vuoto, non monodimensionale o contiene valori non finiti.
\end{apidesc}

\begin{minted}{python}
x = [3.1, 2.7, 4.0, 2.9, 3.5]
print(median(x))  # 3.1
\end{minted}

\section{\texorpdfstring{weighted\_mean(x, w=None)}{weighted\_mean(x, w=None)}}
\label{sec:weighted-mean}
\begin{apidesc}
  \item[\textbf{Scopo}]
    Calcolare la media aritmetica o, se \texttt{w} è fornito, la media pesata
    \[
      \bar{x}_w = \frac{\sum_i w_i x_i}{\sum_i w_i}.
    \]
  \item[\textbf{Restituisce}]
    \texttt{float}: media semplice o pesata di \texttt{x}.
  \item[\textbf{Eccezioni}]
    \texttt{ValueError} se \texttt{x} è vuoto o non finito;
    \texttt{ValueError} se \texttt{w} ha forma incompatibile, contiene valori non positivi o non finiti.
\end{apidesc}

\section{\texorpdfstring{variance(x, w=None, ddof=0)}{variance(x, w=None, ddof=0)}}
\label{sec:variance}
\begin{apidesc}
  \item[\textbf{Scopo}]
    Calcolare la varianza non pesata oppure pesata:
    \[
      \mathrm{Var}_w(x) = \frac{\sum_i w_i (x_i - \bar{x}_w)^2}{\sum_i w_i - \texttt{ddof}}.
    \]
  \item[\textbf{Restituisce}]
    \texttt{float}: varianza di \texttt{x}.
  \item[\textbf{Eccezioni}]
    \texttt{ValueError} se il denominatore è non positivo;
    eredita i controlli di validazione comuni su dati e pesi.
\end{apidesc}

\begin{minted}{python}
x = [1.0, 2.0, 3.0, 4.0, 5.0]
print(variance(x))          # 2.0
print(variance(x, ddof=1))  # 2.5
\end{minted}

\section{\texorpdfstring{covariance(x, y, w=None)}{covariance(x, y, w=None)}}
\label{sec:covariance}
\begin{apidesc}
  \item[\textbf{Scopo}]
    Calcolare la covarianza tramite l'identità
    \[
      \mathrm{Cov}(x,y) = E[xy] - E[x]E[y].
    \]
  \item[\textbf{Restituisce}]
    \texttt{float}: covarianza tra \texttt{x} e \texttt{y}.
  \item[\textbf{Eccezioni}]
    \texttt{ValueError} se \texttt{x} e \texttt{y} non hanno la stessa lunghezza;
    eredita i controlli di validazione comuni su dati e pesi.
\end{apidesc}

\section{\texorpdfstring{standard\_deviation(x, w=None, ddof=0)}{standard\_deviation(x, w=None, ddof=0)}}
\label{sec:standard-deviation}
\begin{apidesc}
  \item[\textbf{Scopo}]
    Calcolare la deviazione standard come radice quadrata della varianza.
  \item[\textbf{Restituisce}]
    \texttt{float}: deviazione standard di \texttt{x}.
  \item[\textbf{Eccezioni}]
    Eredita le stesse eccezioni di \texttt{variance}.
\end{apidesc}

\begin{minted}{python}
x = [2.0, 4.0, 4.0, 4.0, 5.0, 5.0, 7.0, 9.0]
print(standard_deviation(x))          # 2.0
print(standard_deviation(x, ddof=1))  # 2.138...
\end{minted}
